\def\sampleids{SAMPLE_00000,SAMPLE_00001,SAMPLE_00002,SAMPLE_00005,SAMPLE_00008}
\def\samplefigdir{figures/}

\documentclass[11pt,a4paper]{article}

\usepackage[utf8]{inputenc}
\usepackage[T1]{fontenc}
\usepackage{lmodern}
\usepackage[english]{babel}
\usepackage{geometry}
\geometry{a4paper, left=2.5cm, right=2.5cm, top=2.5cm, bottom=2.5cm}
\usepackage{graphicx}
\usepackage{microtype}
\usepackage{amsmath,amssymb,amsfonts}
\usepackage{booktabs}
\usepackage{multirow}
\usepackage{array}
\usepackage{caption}
\usepackage{subcaption}
\usepackage{xcolor}
\usepackage{hyperref}
\usepackage{siunitx}
\usepackage{algorithm}
\usepackage{algpseudocode}

\graphicspath{{figures/}}
\hypersetup{colorlinks=true, linkcolor=blue, citecolor=blue, urlcolor=blue}
\sisetup{detect-all}

\newcommand{\rps}{\ensuremath{\mathrm{RPS}}}
\newcommand{\revps}{\ensuremath{\mathrm{rev}/\mathrm{s}}}
\newcommand{\rsq}{\ensuremath{R^{2}}}

\title{\textbf{Classical baselines for multi-rotor \rps{} estimation}\\
  \large Why they fail and by how much}
\author{Dmitrii Mukhutdinov}
\date{\today}

\begin{document}
\maketitle

% ============================================================================
\section{Classical methods we tried}
% ============================================================================

The task is to recover four per-rotor speeds from a single-channel mixture of
speech + drone noise at very low SNR ($-30$ to $0$~dB).  Before training a
neural network it is worth checking whether standard signal-processing methods
can do the job.  We picked four well-known single-pitch estimators and extended
each to the multi-rotor case with the same greedy suppression procedure.

\paragraph{PYIN.} The probabilistic YIN tracker from \texttt{librosa}.  It is
a single-$f_0$ method; we run it with a wide frequency range $[50, 400]$~Hz to
capture both the motor rotation frequency and the blade-passage frequency, then
divide by the blade count ($b=2$) when necessary.  Because PYIN can only track
one source, we replicate its single output for all four rotors.

\paragraph{Cepstral analysis.} For each STFT frame we compute the real
cepstrum $c = \mathcal{F}^{-1}\{\log|\mathcal{F}\{\mathbf{s}\}|\}$ and search
for the strongest peak in the quefrency band that maps to the plausible
\rps{} range $[50, 150]$~\revps.  The quefrency peak is converted back to
frequency via $f = f_s / q$ and then to rotor speed $\rho = f / b$.

\paragraph{Harmonic Product Spectrum (HPS).} The spectrum is multiplied with
downsampled copies of itself:
\begin{equation}
  P_{\text{HPS}}[k] = |S[k]| \cdot |S[2k]| \cdot |S[3k]| \cdot |S[4k]|.
\end{equation}
The peak of $P_{\text{HPS}}$ in the band $[b\rho_{\min}, b\rho_{\max}]$
gives the blade-passage frequency; dividing by $b$ yields the rotor speed.

\paragraph{Matched-filter bank.} We build a dense grid of candidate rotor
speeds $\rho \in \{50, 50.5, \dots, 150\}$~\revps{} and, for each candidate,
a binary comb template with narrow windows around its first 15 harmonics.  The
template that maximises the cosine similarity with the magnitude spectrum is
taken as the estimate.

\paragraph{NMF with harmonic dictionary.}  Semi-supervised non-negative matrix
factorisation: the magnitude spectrogram $\mathbf{V}$ is factorised as
$\mathbf{V} \approx \mathbf{W}\mathbf{H}$ where $\mathbf{W}$ is a fixed
dictionary of Gaussian harmonic-comb spectral templates (one column per
candidate RPS, same grid as the matched filter).  $\mathbf{H}$ is learned via
multiplicative updates; for each frame the four strongest activations are
taken as the rotor-speed estimates.  This is a true multi-pitch method---there
is no greedy suppression step.

All methods share the same constants: $f_s = \SI{16}{\kilo\hertz}$,
$N=2048$, $H=512$, blade count $b=2$, and 15 harmonics.

% ============================================================================
\section{The greedy multi-rotor extension}
% ============================================================================

The first four methods estimate \emph{one} rotor speed per frame.  To get
four estimates we use the same greedy suppression scheme for all of them:

\begin{algorithm}[H]
\caption{Greedy multi-rotor extractor}
\label{alg:greedy}
\begin{algorithmic}[1]
\Require Frame spectra $\mathbf{S} \in \mathbb{R}^{T \times F}$,
  single-rotor estimator $\mathcal{E}$,
  rotors $R=4$,
  suppression half-width $w=3$ bins
\Ensure $\hat{\mathbf{R}} \in \mathbb{R}^{R \times T}$
\For{$t = 1 \dots T$}
  \State $\mathbf{s} \gets \mathbf{S}[t,:]$\qquad\Comment{copy}
  \For{$r = 1 \dots R$}
    \State $\hat{\rho}_{r,t} \gets \mathcal{E}(\mathbf{s})$
    \For{$h = 1 \dots 15$}
      \State $k_h \gets \text{round}(h \, b \, \hat{\rho}_{r,t} \, N / f_s)$
      \State $\mathbf{s}[k_h-w \;\dots\; k_h+w] \gets 0$
    \EndFor
  \EndFor
\EndFor
\end{algorithmic}
\end{algorithm}

This is simple, deterministic, and wrong in several ways.  Hard zeroing removes
not only the target rotor but also overlapping speech energy and sidebands from
other rotors, so the spectrum gets progressively more damaged for the second,
third and fourth rotors.  Each frame is processed independently, so there is no
temporal continuity.  And the first rotor always gets the cleanest spectrum while
the others fight for scraps.  These are not implementation bugs; they are
structural limitations that any single-pitch method extended this way will hit.

% ============================================================================
\section{Learned baseline for comparison}
% ============================================================================

SimpleConv is a small CNN ($\approx$1~M parameters) that maps the
log-magnitude STFT of the mixture directly to four per-rotor speed traces via
frame-wise regression.  It has no explicit harmonic model, no suppression step,
and no hand-designed peak picking.  The encoder is five 2-D conv blocks that
downsample frequency while keeping time intact; a spatial mean pools the
remaining frequency bins and a two-layer 1-D conv head outputs four channels.  It
is trained end-to-end on DREGON-LM by MSE against telemetry-aligned ground
truth.  See the architecture table in the main paper if you care about the exact
layer dimensions; for this report the only thing that matters is that it is a
plain neural network that learns the mapping from data.

% ============================================================================
\section{Results}
% ============================================================================

\subsection{Numbers}

Table~\ref{tab:summary} shows mean and standard deviation across the 10
DREGON-LM test samples.  The classical methods are not just worse---they are
\emph{two to three orders of magnitude} worse.

\begin{table}[ht]
  \centering
  \caption{Aggregate metrics on 10 test samples.}
  \label{tab:summary}
  \small
  \begin{tabular}{lccc}
\toprule
 & MSE & MAE & $R^2$ \\
method &  &  &  \\
\midrule
SimpleConv & 8.44 ± 14.35 & 1.76 ± 1.17 & 0.7757 ± 0.1506 \\
PYIN & 630.75 ± 602.91 & 22.14 ± 8.01 & -27.0371 ± 14.4139 \\
Cepstral & 1383.49 ± 503.47 & 28.85 ± 4.77 & -79.4142 ± 38.5403 \\
HPS & 868.48 ± 139.21 & 27.54 ± 1.35 & -53.7031 ± 26.2931 \\
Matched Filter & 1343.50 ± 296.38 & 29.83 ± 2.84 & -80.8404 ± 37.4684 \\
NMF & 617.17 ± 430.70 & 18.95 ± 5.04 & -31.0281 ± 18.0616 \\
\bottomrule
\end{tabular}

\end{table}

NMF is the best classical method (MSE~$615$), roughly $86\times$ worse than
SimpleConv but clearly ahead of the greedy single-pitch methods
($870$--$1380$).  Cepstral and matched-filter are $193\times$ and
$187\times$ worse.  The \rsq{} values are deep negative, which means even
the best classical predictions are worse than predicting the per-sample mean.
In other words, these methods do not just fail to track the rotors; they
actively make things worse.

\begin{figure}[ht]
  \centering
  \includegraphics[width=0.85\textwidth]{fig_bar_mse}
  \caption{Mean MSE on a log scale.  Error bars are $\pm$1 std.}
  \label{fig:bar}
\end{figure}

Table~\ref{tab:per_sample} breaks it down per sample.  SimpleConv stays below
$40$ on every single clip.  The classical methods never drop below $160$ and
routinely exceed $1000$.

\begin{table}[ht]
  \centering
  \caption{Per-sample MSE [(rev/s)$^2$].}
  \label{tab:per_sample}
  \small
  \begin{tabular}{lcccccc}
\toprule
 & SimpleConv & PYIN & Cepstral & HPS & Matched Filter & NMF \\
sample &  &  &  &  &  &  \\
\midrule
 & NaN & NaN & NaN & NaN & NaN & NaN \\
 & NaN & NaN & NaN & NaN & NaN & NaN \\
 & NaN & NaN & NaN & NaN & NaN & NaN \\
 & NaN & NaN & NaN & NaN & NaN & NaN \\
 & NaN & NaN & NaN & NaN & NaN & NaN \\
 & NaN & NaN & NaN & NaN & NaN & NaN \\
 & NaN & NaN & NaN & NaN & NaN & NaN \\
 & NaN & NaN & NaN & NaN & NaN & NaN \\
 & NaN & NaN & NaN & NaN & NaN & NaN \\
 & NaN & NaN & NaN & NaN & NaN & NaN \\
\bottomrule
\end{tabular}

\end{table}

\subsection{What the predictions actually look like}

Mean RPS is misleading because the four rotors move independently and the mean
can be flat even when individual rotors are wrong.  Figures~\ref{fig:page00}
to~\ref{fig:page05} show the per-rotor predictions for individual samples.
Each page is one sample: spectrogram on top, then six per-rotor panels (one
per method).  Dotted lines are ground truth; solid lines are predictions.
Rotor colours are consistent across all panels.

\begin{figure}[p]
  \centering
  \includegraphics[width=\textwidth]{fig_page_sample_00000}
  \caption{Sample~00000 -- the easy case.  SimpleConv tracks all four rotors
    closely; the classical methods are essentially flat or random.}
  \label{fig:page00}
\end{figure}

\begin{figure}[p]
  \centering
  \includegraphics[width=\textwidth]{fig_page_sample_00002}
  \caption{Sample~00002 -- a harder clip where SimpleConv MSE rises to $38$
    but still remains two orders of magnitude below the classical methods.}
  \label{fig:page02}
\end{figure}

\begin{figure}[p]
  \centering
  \includegraphics[width=\textwidth]{fig_page_sample_00005}
  \caption{Sample~00005 -- one of the few clips where PYIN finds voiced frames.
    It still tracks only one periodicity and replicates it for all four rotors.}
  \label{fig:page05}
\end{figure}

A few things are immediately obvious from the plots:

\begin{itemize}
  \item \textbf{SimpleConv} actually follows the four dotted traces.  It is not
    perfect---there is some overshoot and occasional mixing of rotors---but the
    solid lines are clearly correlated with the ground truth.
  \item \textbf{PYIN} returns a single trace replicated four times.  When it
    finds voiced frames it locks onto whatever is strongest in the mixture,
    which may be speech or a single rotor, and ignores the other three.
  \item \textbf{Cepstral, HPS, and Matched Filter} were at least extended with
    greedy suppression, but the second, third and fourth estimates are driven by
    random spectral peaks because hard bin zeroing is too destructive.
  \item \textbf{NMF} (bottom panel) does better than the greedy methods: its
    traces sometimes follow the correct rotor for a few frames before losing
    track.  But it still drifts and confuses rotors when harmonics overlap.
\end{itemize}

% ============================================================================
\section{Why they fail}
% ============================================================================

The main reason is embarrassingly obvious: four of the five methods are not
multi-pitch estimators.

Cepstral analysis, HPS, and the matched-filter bank are all single-pitch
algorithms.  They are designed to find one dominant periodicity in a spectrum.
To make them handle four rotors we added a greedy extension that finds the
strongest pitch, suppresses its harmonics by zeroing out bins, and repeats.
This is not how multi-pitch estimation works in the literature; it is a crude
hack that we tried because it was easy to implement.  The problem is that
zeroing a few bins around a harmonic is highly destructive: it removes not only
the target rotor but also overlapping speech formants, sidebands from other
rotors, and any rotor energy that has leaked beyond the narrow suppression
window.  By the time we search for the third or fourth rotor, the spectrum is
so mangled that the estimate is essentially noise.

PYIN is even worse off because we did not even attempt the greedy adaptation.
It is a single-$f_0$ tracker; we just run it once per frame and replicate the
single output four times.  Of course it cannot resolve four independently
varying rotors.

NMF is different: it is a true multi-pitch method.  It jointly decomposes the
spectrum into four harmonic combs without hard suppression.  Yet it still
fails, albeit by a smaller margin (MSE~$615$ vs $870$--$1380$ for the greedy
methods).  Why?  Overlapping harmonics from four rotors create ambiguity in
the dictionary: many atoms have similar harmonic patterns, so the multiplicative
updates spread energy across neighbouring candidates.  Speech adds non-harmonic
energy that NMF tries to explain with the same harmonic dictionary, creating
spurious activations.  And without a temporal model, the frame-wise
factorisation can flip which activation is assigned to which rotor from one
frame to the next.  NMF improves on the greedy single-pitch hack, but the
multi-pitch + speech problem is still too hard for a classical spectral
decomposition with no learned temporal structure.

% ============================================================================
\section{Sanity check: what happens without speech and with only one rotor?}
% ============================================================================

So far the classical methods look useless.  But remember they are
single-pitch estimators, and DREGON-LM is a multi-pitch + speech mixture.
Before writing them off completely, we checked how they behave on the original
DREGON individual-motor recordings: clean recordings of a single rotor
running at constant speed (no speech, no other rotors).  We trimmed the start
and end transients and evaluated the \emph{single-pitch} versions of each
method---no greedy suppression, no multi-rotor hack.

The results are summarised in Table~\ref{tab:single_rotor} and
Figure~\ref{fig:single_rotor_scatter}.  PYIN and the matched filter are
essentially perfect on most of these clips (MAE below $2$~\revps{} when they
detect a periodicity).  Even cepstral analysis and HPS, while noisier, stay
within $5$--$10$~\revps{} of the target.  NMF is mixed: on some motors it is
accurate to within a few rev/s, but on others it locks onto the wrong harmonic
or a spurious spectral peak and drifts far from the target.  The single-pitch
methods are not broken; they are simply solving the wrong problem on DREGON-LM.

\begin{table}[ht]
  \centering
  \caption{Mean MSE / MAE on clean single-rotor recordings (single-pitch
    methods, no greedy suppression).  SimpleConv is shown as the best rotor
    and averaged over all four outputs.}
  \label{tab:single_rotor}
  \small
  \begin{tabular}{lcc}
\toprule
Method & Mean MSE & Mean MAE \\
\midrule
PYIN & 240.8 & 6.2 \\
Cepstral & 742.7 & 15.1 \\
HPS & 1149.5 & 20.7 \\
Matched Filter & 168.2 & 6.9 \\
NMF & 599.7 & 11.5 \\
SimpleConv (best rotor) & 2027.3 & 41.1 \\
SimpleConv (avg over 4) & 2078.2 & 42.0 \\
\bottomrule
\end{tabular}

\end{table}

\begin{figure}[ht]
  \centering
  \includegraphics[width=0.7\textwidth]{fig_single_rotor_scatter}
  \caption{Mean predicted RPS vs target RPS on clean single-rotor recordings.
    Classical single-pitch methods cluster around the diagonal; SimpleConv
    predicts random low values because it has never seen a single-rotor input.}
  \label{fig:single_rotor_scatter}
\end{figure}

The really interesting result, however, is SimpleConv.  On DREGON-LM it
tracks all four rotor speeds quite accurately, so it is not simply
memorising chunks of training audio; it generalises to unseen mixtures.
On single-rotor recordings, though, it fails: all four channels predict
roughly the same low value ($10$--$50$~\revps{} instead of the true
$50$--$90$), so the four dotted lines in
Figures~\ref{fig:single_rotor_set1} and \ref{fig:single_rotor_set2}
overlap perfectly.  The network has clearly learned that there should be
four rotors and that their speeds should differ and vary over time.  When
presented with a single constant rotor, it cannot reconcile that with its
internal model of the problem, and all four channels collapse to the same
confused estimate.

\begin{figure}[p]
  \centering
  \includegraphics[width=\textwidth]{fig_single_rotor_set1}
  \caption{Predictions on clean single-rotor recordings (Motors~1 at 50, 70,
    90~\revps).  Classical single-pitch methods track the constant target well.
    SimpleConv outputs four channels (dotted lines), but on single-rotor input
    all four collapse to the same wrong value, so the four dotted lines overlap
    perfectly.  The solid line is their mean (identical to each channel), and
    one channel is shown bold to mark the ``best'' channel (they are all equal).}
  \label{fig:single_rotor_set1}
\end{figure}

\begin{figure}[p]
  \centering
  \includegraphics[width=\textwidth]{fig_single_rotor_set2}
  \caption{Predictions on clean single-rotor recordings (Motors~4 at 50, 70,
    90~\revps).  Same pattern as Figure~\ref{fig:single_rotor_set1}: all four
    SimpleConv channels output the same wrong value, so the dotted lines overlap.}
  \label{fig:single_rotor_set2}
\end{figure}

\begin{figure}[p]
  \centering
  \includegraphics[width=\textwidth]{fig_single_rotor_allmotors}
  \caption{Predictions on the four-rotor synchronized recording (allMotors\_70).
    Classical methods struggle with multiple overlapping combs.  SimpleConv
    dotted = four channels, solid = mean, bold = best channel; the four outputs
    are now clustered near the true speed because the input matches training.}
  \label{fig:single_rotor_allmotors}
\end{figure}

On the four-rotor synchronized recording (allMotors\_70,
Figure~\ref{fig:single_rotor_allmotors}) SimpleConv performs much better
($\approx 75$--$83$~\revps{} vs target $70$).  The input now contains four
rotors, which matches the network's structural expectation; the absence of
speech is apparently less important than the presence of the expected number
of rotors.

% ============================================================================
\section{Conclusion}
% ============================================================================

On the DREGON-LM test set (multi-pitch + speech, low SNR) the classical
baselines fail by two to three orders of magnitude versus a 1-M-parameter
CNN.  The four single-pitch methods (PYIN, cepstral, HPS, matched filter)
are forced into a multi-pitch task with a crude greedy suppression hack that
wipes out useful information.  NMF, a true multi-pitch method, improves on
the greedy approaches (MSE~$615$ vs $870$--$1380$) but is still
$86\times$ worse than SimpleConv and produces drifting, confused traces.

On clean single-rotor recordings the same classical methods are accurate to
within a few rev/s.  They are not broken; the multi-pitch task is simply
harder than they were built for.

SimpleConv shows the reverse.  It tracks the four rotor speeds accurately
on DREGON-LM, so it is doing more than memorising training chunks; it
generalises to new mixtures.  Yet on single-rotor recordings all four
channels collapse to the same wrong number ($10$--$50$~\revps{} instead of
$50$--$90$).  The fact that they all emit the same value is telling: the
network has learned that there should be four rotors with different,
varying speeds, and when that structure is absent it cannot produce a
meaningful estimate.  On the four-rotor recording (allMotors\_70) it
recovers speeds close to the target ($\approx 75$--$83$ vs $70$), because
the input now matches the expected four-rotor structure.

Bottom line: even a proper multi-pitch classical method (NMF with harmonic
priors) cannot solve this task reliably at low SNR.  SimpleConv succeeds
because it has learned the specific structure of the problem (four rotors,
variable speeds, mixed with speech).  It is not a general rotor-speed
estimator, and its single-rotor failure shows it is relying on structural
expectations rather than on independent per-rotor harmonic analysis.

\end{document}
